\section{Besov范数}
\label{sec:besov-norms}

上一节得到了两类信息：连续模$\omega_r(f,t)_p$记录尺度$t$上的变化，积分指数$p$决定同一尺度内怎样度量位置不同的变化；不同尺度之间怎样汇总则尚未规定。本节把小尺度变化乘以光滑性权重$t^{-s}$，再对所有尺度统一汇总，定义Besov范数。随后用周期阶跃检验每个参数的作用，最后证明空间完备。

\subsection{差分范数}
\label{sec:besov-difference-norm}

汇总不同尺度时，要让每个尺度得到同样的分量。测度$\mathrm dt/t$满足这一点：每个二进区间$[2^{-j-1},2^{-j}]$的测度都是$\ln2$，与$j$无关。在这个测度下取$L^q$大小，就是把各尺度的贡献按$\ell^q$方式汇总。

\begin{definition}[周期Besov空间]
\label{def:besov-space}
设$1\leq p,q\leq\infty$，$s>0$，整数$r>s$。对$f\in L^p(\mathbb T)$，令
\begin{equation*}
\begin{aligned}
|f|_{B^s_{p,q}}&=\Bigl(\int_0^1\bigl(t^{-s}\omega_r(f,t)_p\bigr)^q\,\frac{\mathrm dt}{t}\Bigr)^{1/q}\quad(q<\infty),\\
|f|_{B^s_{p,\infty}}&=\sup_{0<t\leq1}t^{-s}\omega_r(f,t)_p ,
\end{aligned}
\end{equation*}
并令$\lVert f\rVert_{B^s_{p,q}}=\lVert f\rVert_p+|f|_{B^s_{p,q}}$。所有使$\lVert f\rVert_{B^s_{p,q}}<\infty$的$f\in L^p(\mathbb T)$组成的空间称为周期\term{Besov空间}（Besov space），记为$B^s_{p,q}(\mathbb T)$，简记为$B^s_{p,q}$。
\end{definition}

定义中$t^{-s}\omega_r(f,t)_p$是尺度$t$上的变化相对于期望量级$t^s$的比值，$|f|_{B^s_{p,q}}$把这些比值按$\ell^q$方式汇总。只取$t\leq1$是因为周期为$1$，更大的步长不提供新信息；大尺度的大小由$\lVert f\rVert_p$记录。不同的$r>s$给出等价的范数，所以空间不依赖$r$的选择，这一点这里不给出证明。

由于$\omega_r$单调，积分形式可以换成离散的二进和。对$t\in[2^{-j-1},2^{-j}]$，$\omega_r(f,2^{-j-1})_p\leq\omega_r(f,t)_p\leq\omega_r(f,2^{-j})_p$，且$2^{js}\leq t^{-s}\leq2^{(j+1)s}$，每段的$\mathrm dt/t$测度为$\ln2$。逐段估计后求和，得
\begin{equation}
2^{-sq}\ln2\sum_{j\geq1}\bigl(2^{js}\omega_r(f,2^{-j})_p\bigr)^q\leq|f|_{B^s_{p,q}}^q\leq2^{sq}\ln2\sum_{j\geq0}\bigl(2^{js}\omega_r(f,2^{-j})_p\bigr)^q ,
\label{eq:besov-dyadic-form}
\end{equation}
$q=\infty$时把和换成上确界。$j=0$的项不超过$2^r\lVert f\rVert_p$，已被范数的$L^p$部分控制。所以$|f|_{B^s_{p,q}}$与序列$(2^{js}\omega_r(f,2^{-j})_p)_{j\geq0}$的$\ell^q$范数相差有界倍数。

离散形式立即给出两类包含关系。若$q_1\leq q_2$，由$\ell^{q_1}\subset\ell^{q_2}$，$B^s_{p,q_1}\subset B^s_{p,q_2}$：汇总指数越小，要求越严。若$s'<s$，对任意$q$，$(2^{js'}\omega_r)^q=2^{-j(s-s')q}(2^{js}\omega_r)^q$，后一因子有界时前者可和，所以$B^s_{p,\infty}\subset B^{s'}_{p,q}$：降低一点光滑阶数，就可以换取任意的汇总方式。改变$q$只改变允许的尺度分布，改变$s$则改变每个尺度上期望的量级。

\subsection{参数的作用}
\label{sec:besov-parameters}

三个参数分工明确：$s$规定期望的变化阶数，$p$汇总同一尺度内不同位置的变化，$q$汇总不同尺度的变化。我们用周期阶跃$f$检验这一分工。固定$p=2$与$0<s<1$，此时可以取$r=1$。

由式\eqref{eq:besov-step-difference}，$0<t\leq1/2$时$\omega_1(f,t)_2=(2t)^{1/2}$。$1/2<t\leq1$时，取$h=1/2$得到$\lVert\Delta_hf\rVert_2=1$；而$1/2<h\leq1$时，周期性给出$\Delta_hf=\Delta_{h-1}f$，其范数$(2(1-h))^{1/2}\leq1$，所以$\omega_1(f,t)_2=1$。代入定义，$[1/2,1]$上的部分总是有限，关键是$t\to0$的部分：
\[
\int_0^{1/2}\bigl(t^{-s}(2t)^{1/2}\bigr)^q\,\frac{\mathrm dt}{t}=2^{q/2}\int_0^{1/2}t^{q(1/2-s)-1}\,\mathrm dt .
\]

当$s<1/2$时，指数$q(1/2-s)-1>-1$，积分对每个有限$q$都有限，$q=\infty$时上确界也有限。所以阶跃函数属于每个$B^s_{2,q}$，$0<s<1/2$。当$s=1/2$时，被积函数为$2^{q/2}t^{-1}$，积分对每个有限$q$都发散；而$q=\infty$时，$\sup_tt^{-1/2}\omega_1(f,t)_2=\sqrt2$有限。所以阶跃函数属于$B^{1/2}_{2,\infty}$，不属于任何$q<\infty$的$B^{1/2}_{2,q}$。当$s>1/2$时，$t^{-s}(2t)^{1/2}\to\infty$，阶跃函数不属于任何$B^s_{2,q}$。

这三种情形看清了$s$与$q$各自的作用。$s$决定比较的基准$t^s$：$s<1/2$时阶跃函数在每个尺度上的变化都比基准小得多，并且越细越小，任何汇总方式都可以接受；$s>1/2$时变化比基准大，无论怎样汇总都不行。$s=1/2$恰好是阶跃函数的变化阶数，此时每个二进尺度贡献大小相同的比值$\sqrt2$，只有不要求跨尺度可和的$q=\infty$能够接受。上一节的$f_b$也是如此：在$B^s_{2,q}$中的成员资格大致取决于权重序列$(b_j)$是否属于$\ell^q$。

积分指数$p$的作用也可以从同一个函数读出。由式\eqref{eq:besov-step-difference}，$\omega_1(f,t)_p=(2t)^{1/p}$，所以同样的论证说明阶跃函数属于$B^{1/p}_{p,\infty}$，并属于每个$s<1/p$的$B^s_{p,q}$。$p=1$时阶数为$1$，$p$越大阶数越低：变化集中在小集合上，积分指数越高，这种集中就被看得越重。对比之下，阶跃函数不属于任何$W^{1,p}$；Besov尺度给了它一个确切的位置。

\subsection{完备性}
\label{sec:besov-completeness}

先核对$\lVert\cdot\rVert_{B^s_{p,q}}$是范数。齐次性显然。差分对函数线性，由上一节$\omega_r(f+g,t)_p\leq\omega_r(f,t)_p+\omega_r(g,t)_p$，再对$L^q(\mathrm dt/t)$用Minkowski不等式（$q=\infty$时直接取上确界），得$|f+g|_{B^s_{p,q}}\leq|f|_{B^s_{p,q}}+|g|_{B^s_{p,q}}$。所以Besov条件对加法和数乘封闭，$B^s_{p,q}$是线性空间。若$\lVert f\rVert_{B^s_{p,q}}=0$，则$\lVert f\rVert_p=0$，$f$是零等价类；单独的$|\cdot|_{B^s_{p,q}}$只是半范数，在常数上为零，$L^p$项正是为排除这一点而加入的。

\begin{theorem}[Besov空间的完备性]
\label{thm:besov-complete}
设$1\leq p,q\leq\infty$，$s>0$，整数$r>s$。则$B^s_{p,q}(\mathbb T)$是Banach空间。
\end{theorem}

\begin{proof}
设$(f_n)$是$B^s_{p,q}$中的Cauchy列。由$\lVert f_n-f_m\rVert_p\leq\lVert f_n-f_m\rVert_{B^s_{p,q}}$，它也是$L^p$中的Cauchy列，由$L^p$的完备性收敛到某个$f\in L^p(\mathbb T)$。

下面控制$f-f_m$的差分。固定$m$与$h$。由$\lVert\Delta_h^rg\rVert_p\leq2^r\lVert g\rVert_p$，差分在$L^p$中连续，所以$\Delta_h^r(f_n-f_m)\to\Delta_h^r(f-f_m)$，从而
\begin{equation*}
\begin{aligned}
\lVert\Delta_h^r(f-f_m)\rVert_p&=\lim_{n\to\infty}\lVert\Delta_h^r(f_n-f_m)\rVert_p\\
&\leq\liminf_{n\to\infty}\omega_r(f_n-f_m,t)_p,\qquad|h|\leq t .
\end{aligned}
\end{equation*}
对$|h|\leq t$取上确界，得$\omega_r(f-f_m,t)_p\leq\liminf_n\omega_r(f_n-f_m,t)_p$。乘以$t^{-s}$后，$q<\infty$时由Fatou引理，$q=\infty$时直接取上确界，得
\[
|f-f_m|_{B^s_{p,q}}\leq\liminf_{n\to\infty}|f_n-f_m|_{B^s_{p,q}} .
\]

任给$\varepsilon>0$，取$M$使$m,n\geq M$时$\lVert f_n-f_m\rVert_{B^s_{p,q}}<\varepsilon$。则$m\geq M$时$|f-f_m|_{B^s_{p,q}}\leq\varepsilon$，并且$\lVert f-f_m\rVert_p=\lim_n\lVert f_n-f_m\rVert_p\leq\varepsilon$。于是$f=f_m+(f-f_m)\in B^s_{p,q}$，且$\lVert f-f_m\rVert_{B^s_{p,q}}\leq2\varepsilon$，即$f_m\to f$在$B^s_{p,q}$中成立。
\end{proof}

证明的结构与Sobolev空间的完备性相似：先在较弱的$L^p$中取得极限，再说明附加的条件在极限中保留。不同之处在于，这里附加的条件不是测试恒等式，而是对每个步长的差分估计，它通过差分在$L^p$中的连续性与Fatou引理传到极限。
